
\documentclass[12pt,a4paper]{article}

\usepackage[utf8]{inputenc}
\usepackage[T1]{fontenc}
\usepackage{lmodern}
\usepackage{amsmath}
\usepackage{amssymb}
\usepackage{geometry}
\usepackage{enumitem}

\geometry{margin=2.5cm}

\setlength{\parindent}{0pt}
\setlength{\parskip}{8pt}

\title{\textbf{Tema 4 - Retele Markov: independenta conditionata si factorizare}}
\author{}
\date{}

\begin{document}

\maketitle

% =========================================================
\section{Formulele}
% =========================================================

\subsection{Probabilitatea conditionata}

Formula:

\[
\boxed{
P(A\mid B)=\frac{P(A,B)}{P(B)}
}
\]

\subsection{Regula lantului}

Pentru trei variabile aleatoare:

\[
\boxed{
P(a,b,c)=P(a\mid b,c)P(c\mid b)P(b)
}
\]

\subsection{Independenta conditionata}

\[
\boxed{A\perp C\mid B}
\]

\[
\boxed{
P(a\mid b,c)=P(a\mid b)
}
\]

O forma echivalenta este:

\[
\boxed{
P(a,c\mid b)=P(a\mid b)P(c\mid b)
}
\]

\textbf{Interpretare:}

Daca B este cunoscut, informatia despre C nu mai modifica probabilitatea lui A.

\subsection{Factorizarea unei retele Markov}

Conform teoremei Hammersley-Clifford, pentru o distributie strict pozitiva care respecta proprietatile Markov ale grafului:

\[
\boxed{
P(x)=\frac{1}{Z}
\prod_{C\in\mathcal{C}}\psi_C(x_C)
}
\]

Unde:

\begin{itemize}
    \item $P(x)$ = probabilitatea configuratiei complete.
    \item $\psi_C(x_C)$ = potentialul asociat unei clici.
    \item $\mathcal{C}$ = familia clicilor folosite in factorizare.
    \item $Z$ = constanta de normalizare.
\end{itemize}

\subsection{Constanta de normalizare}

\[
\boxed{
Z=\sum_x\prod_{C\in\mathcal{C}}\psi_C(x_C)
}
\]

\textbf{Ce face Z?}

Asigura faptul ca suma probabilitatilor tuturor configuratiilor este 1.

\[
\sum_x P(x)=1
\]

% =========================================================
\section{Exercitiul 1 - Enunt}
% =========================================================

Consideram trei variabile aleatoare discrete $A$, $B$, $C$ si graful neorientat:

\[
\boxed{A-B-C}
\]

Presupunem ca distributia comuna este strict pozitiva:

\[
P(a,b,c)>0
\]

pentru orice configuratie posibila.

Trebuie sa demonstram:

\textbf{a)} Semnificatia independentei conditionate:

\[
A\perp C\mid B
\]

\textbf{b)} Independenta conditionata implica factorizarea:

\[
P(a,b,c)=
\frac{1}{Z}\phi_{AB}(a,b)\phi_{BC}(b,c)
\]

\textbf{c)} Factorizarea implica independenta conditionata.

% =========================================================
\section{Rezolvare - Punctul a)}
% =========================================================

\subsection*{Cerinta}

Explicati ce inseamna:

\[
A\perp C\mid B
\]

Ce informatie suplimentara despre A mai ofera C atunci cand cunoastem B?

\subsection*{Explicatie intuitiva}

Consideram trei statii meteorologice:

\begin{itemize}
    \item A = statia meteo din Iasi.
    \item B = statia meteo din Vaslui.
    \item C = statia meteo din Barlad.
\end{itemize}

Modelul presupune ca informatia dintre A si C trece prin B.

\[
\underbrace{\text{Iasi}}_A
-
\underbrace{\text{Vaslui}}_B
-
\underbrace{\text{Barlad}}_C
\]

Sa presupunem ca stim deja ca la Vaslui ploua.

Daca aflam si ca la Barlad ploua, aceasta informatie nu mai modifica probabilitatea calculata pentru vremea din Iasi, in modelul care respecta independenta conditionata data.

Matematic:

\[
\boxed{
P(a\mid b,c)=P(a\mid b)
}
\]

% =========================================================
\section{Rezolvare - Punctul b)}
% =========================================================

\subsection*{Cerinta}

Demonstrati ca independenta conditionata implica existenta potentialelor strict pozitive $\phi_{AB}$, $\phi_{BC}$ si a unei constante $Z>0$, astfel incat:

\[
\boxed{
P(a,b,c)=
\frac{1}{Z}\phi_{AB}(a,b)\phi_{BC}(b,c)
}
\]

\subsection*{Pasul 1. Pornim de la regula lantului}

Formula generala:

\[
P(a,b,c)=P(a\mid b,c)P(c\mid b)P(b)
\]

Aceasta formula rezulta prin aplicarea succesiva a definitiei probabilitatii conditionate.

\subsection*{Pasul 2. Aplicam independenta conditionata}

Din punctul a) avem:

\[
A\perp C\mid B
\]

Asadar:

\[
P(a\mid b,c)=P(a\mid b)
\]

Inlocuim in formula:

\[
\boxed{
P(a,b,c)=P(a\mid b)P(c\mid b)P(b)
}
\]

\subsection*{Pasul 3. Definim potentialele}

Vrem sa obtinem un produs intre un factor care depinde numai de A si B si un factor care depinde numai de B si C.

Alegem:

\[
\boxed{
\phi_{AB}(a,b)=P(a\mid b)P(b)
}
\]

si:

\[
\boxed{
\phi_{BC}(b,c)=P(c\mid b)
}
\]

Inlocuind:

\[
P(a,b,c)=\phi_{AB}(a,b)\phi_{BC}(b,c)
\]

\subsection*{Pasul 4. Alegem constanta Z}

Deoarece produsul este deja o distributie comuna normalizata, putem alege:

\[
\boxed{Z=1}
\]

Prin urmare:

\[
\boxed{
P(a,b,c)=
\frac{1}{Z}\phi_{AB}(a,b)\phi_{BC}(b,c)
}
\]

Deoarece $P(a,b,c)>0$ pentru toate configuratiile, marginalele si probabilitatile conditionate utilizate sunt pozitive.

Deci potentialele definite sunt strict pozitive.

\textbf{Demonstratia este completa.}

\subsection*{Exemplu numeric pentru punctul b)}

Presupunem ca pentru anumite valori $a,b,c$:

\[
P(b)=0.5
\]

\[
P(a\mid b)=0.8
\]

\[
P(c\mid b)=0.6
\]

Calculam:

\[
P(a,b,c)=0.8\cdot0.6\cdot0.5
\]

\[
\boxed{P(a,b,c)=0.24}
\]

Acum calculam potentialele:

\[
\phi_{AB}(a,b)=0.8\cdot0.5=0.4
\]

\[
\phi_{BC}(b,c)=0.6
\]

Verificam:

\[
\phi_{AB}(a,b)\phi_{BC}(b,c)
=
0.4\cdot0.6
\]

\[
\boxed{0.24=0.24}
\]

% =========================================================
\section{Rezolvare - Punctul c)}
% =========================================================

\subsection*{Cerinta}

Demonstrati reciproca.

Daca distributia se factorizeaza astfel:

\[
P(a,b,c)=
\frac{1}{Z}\phi_{AB}(a,b)\phi_{BC}(b,c)
\]

atunci:

\[
\boxed{A\perp C\mid B}
\]

\subsection*{Pasul 1. Scriem definitia probabilitatii conditionate}

Vrem sa calculam:

\[
P(a,c\mid b)
\]

Din definitie:

\[
\boxed{
P(a,c\mid b)=\frac{P(a,b,c)}{P(b)}
}
\]

\subsection*{Pasul 2. Calculam probabilitatea lui B}

Probabilitatea lui B se obtine prin insumarea tuturor valorilor posibile ale lui A si C.

\[
P(b)=\sum_{a'}\sum_{c'}P(a',b,c')
\]

Inlocuim factorizarea:

\[
P(b)=
\frac{1}{Z}
\sum_{a'}\sum_{c'}
\phi_{AB}(a',b)\phi_{BC}(b,c')
\]

Observam ca:

\begin{itemize}
    \item $\phi_{AB}$ nu depinde de $c'$.
    \item $\phi_{BC}$ nu depinde de $a'$.
\end{itemize}

Prin urmare, putem separa sumele:

\[
\boxed{
P(b)=\frac{1}{Z}
\left(
\sum_{a'}\phi_{AB}(a',b)
\right)
\left(
\sum_{c'}\phi_{BC}(b,c')
\right)
}
\]

\subsection*{Pasul 3. Calculam probabilitatea conditionata}

Inlocuim numaratorul si numitorul in formula probabilitatii conditionate:

\[
P(a,c\mid b)=
\frac{
\frac{1}{Z}\phi_{AB}(a,b)\phi_{BC}(b,c)
}{
\frac{1}{Z}
\left(\sum_{a'}\phi_{AB}(a',b)\right)
\left(\sum_{c'}\phi_{BC}(b,c')\right)
}
\]

Factorul $\frac{1}{Z}$ se simplifica.

Obtinem:

\[
P(a,c\mid b)=
\frac{
\phi_{AB}(a,b)\phi_{BC}(b,c)
}{
\left(\sum_{a'}\phi_{AB}(a',b)\right)
\left(\sum_{c'}\phi_{BC}(b,c')\right)
}
\]

Separand produsul in doua fractii:

\[
\begin{aligned}
P(a,c\mid b)
&=
\underbrace{
\frac{\phi_{AB}(a,b)}
{\sum_{a'}\phi_{AB}(a',b)}
}_{P(a\mid b)}
\\[8pt]
&\quad\cdot
\underbrace{
\frac{\phi_{BC}(b,c)}
{\sum_{c'}\phi_{BC}(b,c')}
}_{P(c\mid b)}
\end{aligned}
\]

Deci:

\[
\boxed{
P(a,c\mid b)=P(a\mid b)P(c\mid b)
}
\]

Aceasta este exact definitia independentei conditionate.

Rezulta:

\[
\boxed{A\perp C\mid B}
\]

\textbf{Demonstratia reciprocei este completa.}

% =========================================================
\section{Concluzia exercitiului}
% =========================================================

Am demonstrat ambele implicatii.

\textbf{Prima implicatie:}

\[
A\perp C\mid B
\implies
P(a,b,c)=
\frac{1}{Z}\phi_{AB}(a,b)\phi_{BC}(b,c)
\]

\textbf{A doua implicatie:}

\[
P(a,b,c)=
\frac{1}{Z}\phi_{AB}(a,b)\phi_{BC}(b,c)
\implies
A\perp C\mid B
\]

Prin urmare, in conditiile exercitiului:

\[
\boxed{
A\perp C\mid B
\iff
P(a,b,c)=
\frac{1}{Z}\phi_{AB}(a,b)\phi_{BC}(b,c)
}
\]

unde potentialele sunt strict pozitive, iar $Z>0$.

Acest rezultat ilustreaza, pentru graful simplu $A-B-C$, legatura dintre proprietatile Markov si factorizarea distributiei in potentiale locale.

\end{document}
